\anexo{Entrevista 2: Vendedor de Mercado Libre (foco: validación de la solución)}

Segunda de las tres entrevistas semiestructuradas descriptas en la metodología general del Anexo D. Nombre utilizado es un pseudónimo, en línea con el criterio de anonimización ya explicado (Sección~\ref{sec:viabilidad-legal}).

\emph{Entrevistado}: Lucas, vendedor de Mercado Libre, mostrándole el dashboard del sistema con el ranking y el score de tendencia.

\emph{Objetivo}: validar la comprensión y utilidad del dashboard y del score de tendencia mostrando el sistema en funcionamiento.

\begin{enumerate}
	\item \textbf{¿Qué vendés, hace cuánto, y cómo decidís hoy qué stockear?} Respondió que vende tecnología, y que hoy decide con intuición, combinando redes sociales, el ranking de Mercado Libre y el precio del importador. Lo describió como apostar con datos sueltos.
	\item \textbf{Antes de ver la herramienta: si una aplicación te dijera ``este producto está creciendo'', ¿qué necesitarías saber para creerle?} Indicó que necesitaría saber la fuente del dato (si son ventas reales de Mercado Libre o una estimación), el período de crecimiento, el volumen base (aclaró que pasar de 2 a 10 ventas no es lo mismo que pasar de 500 a 1000) y si el crecimiento es consistente o fue solo un pico de un día. Sin esa información, dijo, el número no significa nada para él.
	\item \textbf{(Mostrando el dashboard) ¿Qué entendés que muestra esta pantalla?} Interpretó que se trata de un ranking de productos ordenados según qué tan en tendencia están, con un puntaje y un desglose de cuánto aparece en búsquedas, hace cuánto viene creciendo, su precio y su posición en el catálogo. En sus palabras, resumió que la pantalla muestra ``lo que está creciendo ahora''.
	\item \textbf{El score va de 0 a 100: ¿qué esperarías de un producto con 80 frente a uno con 40?} Explicó que con un puntaje de 80 esperaría un producto que viene subiendo con fuerza desde hace dos o tres semanas, con volumen real y todavía sin saturar, algo que investigaría el mismo día. Con 40, en cambio, lo consideraría dudoso: podría ser un pico raro o algo recién empezando sin saber si se sostiene, así que lo miraría, pero no arriesgaría stock en base a eso.
	\item \textbf{¿Qué harías mañana con este ranking si fuera tuyo?} Describió que filtraría por su categoría, ordenaría por score y a los tres o cuatro productos mejor rankeados les pediría precio al importador; los que tuvieran poca competencia y buen margen los publicaría esa misma semana, y después configuraría una alerta semanal.
	\item \textbf{De los componentes del score (aparece seguido, se mantiene en el tiempo, posición en el catálogo, precio estable), ¿cuál te importa más? ¿Agregarías alguno?} Dijo que los componentes que más le importan son que se mantenga en el tiempo y la posición en el catálogo, porque aparecer seguido puede ser solo un pico aislado. Como componentes adicionales, propuso sumar la cantidad de vendedores compitiendo (saturación), la velocidad de caída del precio y un margen estimado según el costo típico de la categoría.
	\item \textbf{¿Preferís entrar a mirar el panel o recibir una alerta cuando algo empieza a crecer? ¿Por qué medio?} Prefiere sin dudas la alerta: si tiene que entrar a mirar el panel, lo hace como mucho una vez por semana, y consideró que es justamente la alerta la que evita llegar tarde, que es su problema principal. Sugirió que llegue por WhatsApp o mail, con un mensaje simple del tipo ``este producto está creciendo en tu categoría''.
	\item \textbf{¿Qué le falta a esta herramienta para que la uses en serio en tu operación?} Pidió que el dato sea específico de su categoría y no un ranking genérico de toda la plataforma, que se pueda ver el historial (no solo el score del día) y que se actualice con frecuencia; aclaró que si el dato está desactualizado o es demasiado general, no le sirve para decidir compras de stock.
	\item \textbf{¿La usarías gratis? ¿Pagarías por ella? ¿Qué valor mensual te parecería razonable?} Afirmó que gratis la usaría sin dudarlo, y que también pagaría si de verdad le permitiera entrar dos semanas antes a una tendencia, ya que un solo error de timing en una tendencia le costó más de lo que pagaría en dos años de suscripción. Consideró razonable un rango de entre 10 y 30 dólares mensuales, siempre que la herramienta demuestre que ese costo se recupera en margen.
	\item \textbf{¿Qué es lo que menos te convence de lo que viste?} Lo que menos le convenció fue que el score funcione como una caja negra: si no puede entender por qué un producto tiene 80 puntos, no confía en el dato. También cuestionó que la data de ``apariciones'' no sea la venta real de Mercado Libre, ya que los datos públicos de la plataforma son estimaciones, y pidió que se aclare de dónde sale cada número.
\end{enumerate}
